\par\Needspace{8\baselineskip}\bigskip
\SourceLeaf{487}%
\pdfbookmark[0]{Entwurf zur Gradmessung. 30. Mai 1819.}{entwurf-gradmessung-mai1819}
\begin{center}\large [ENTWURF ZUR GRADMESSUNG.]\end{center}
\par
— — — Über Gradmessungen überhaupt, und in wie fern sie einen der interessantesten Gegenstände des menschlichen Wissens begründen, darf ich in einem Schreiben an Ew. — — nichts sagen. Allein die grossen Vortheile und neuen Aufklärungen noch dunkler Punkte, welche von der Vervielfältigung solcher Operationen zu erwarten sind, beruhen doch mit auf der Bedingung, dass diese so viel als möglich ins Grosse gehen. Isolirte Gradmessungen in Europa, die nur einen kleinen Umfang umfassen, könnten jetzt, nach den grossen Arbeiten in Frankreich und England, nur einen sehr untergeordneten Werth haben: wogegen noch eine oder ein Paar andere ähnliche Messungen in Europa, von einer bedeutenden Ausdehnung, gewiss für die Kenntniss der Gestalt der Erde ungemein wichtig sein würden.

Ein Blick auf die Karte von Europa, und auf den Culturzustand der verschiedenen Nationen, zeigt dass ausser der grossen Linie von den Balearischen bis zu den Orkneys Inseln, nur noch in zwei Richtungen ähnliche Operationen ausführbar sein werden 1) im Russischen Reiche und 2) durch die Jütische Halbinsel und ganz Deutschland bis zum Mittelländischen Meere. Zur Ausführung einer Messung in ersterm Reiche scheinen, mir zugekommenen Nachrichten zufolge, einige Aussichten zu sein. Allein für die andern ist mehr als Aussicht: der erste Hauptschritt ist bereits wirklich geschehen. Die von dem Könige von Dänemark befohlene Gradmessung, von der Nordspitze Jütlands bis Lauenburg ist bereits seit zwei Jahren im Gange; dass dieselbe ganz so wird ausgeführt werden, wie es der heutige Zustand der Wissenschaften und Künste möglich und nothwendig macht, dafür bürgt die Geschicklichkeit und Einsicht des dänischen Astronomen, die Vortrefflichkeit seiner Instrumente, und die Liberalität, womit der König von Dänemark alles, was zur Vollkommenheit der Messung nothwendig oder wünschenswerth gefunden wird, genehmigt.

Die dänische Gradmessung soll, ausser dem erwähnten Meridianbogen, auch noch die Messung des Bogens eines Parallelkreises von der Westküste Jütlands bis Koppenhagen umfassen. Natürlich kann hier nur von dem erstern Theile des Planes die Rede sein. Jener Bogen umfasst für sich schon \(4\frac12\) Grad, und die Messung ist daher, schon isolirt betrachtet, von einer respectabeln Ausdehnung. Allein ohne Vergleich wichtiger erscheint dieselbe, wenn sie als der Anfang jener grossen Messung betrachtet wird, die einer Ausdehnung bis zur Insel Elba, also bis auf \(16^0\), fähig ist. Dass, früh oder spät, diese Operation in einer solchen Ausdehnung einmal werde ausgeführt werden, ist wohl mehr als eine chimärische Hoffnung, besonders da in einigen dazwischen liegenden Ländern, namentlich in Thüringen und Bayern, bereits manche Vorarbeiten wirklich geschehen sind. Die Fortsetzung der dänischen Messung durch das Königreich Hannover ist aber die erste und wesentlichste Bedingung zur dereinstigen Realisirung jenes Planes. Durch diese Fortsetzung allein würde der Bogen schon um 2 Grad vergrössert werden. Würde dann auch noch die Gothaische Sternwarte durch ein Dreiecksnetz mit der Göttingischen in Verbindung gebracht, was auch in mancher andern Rücksicht sehr wünschenswerth und leicht ausführbar sein würde, so wäre dadurch schon ein Bogen von 7 Graden realisirt.

Nur kurz brauche ich zu berühren, dass die Messung eines Meridianbogens von Hamburg bis Göttingen auch noch in andern Beziehungen, als der reinwissenschaftlichen, von grosser Wichtigkeit sein würde. Das zu diesem Zweck geführte Dreiecksnetz würde, wenn über kurz oder lang eine den heutigen Forderungen entsprechende Vermessung des ganzen Königreichs Hannover beschlossen werden sollte, die sicherste Grundlage abgeben, um die weitern Triangulationen östlich und westlich an dasselbe anzuschliessen. Und falls zu einer solchen Generalvermessung nahe Aussicht sein sollte, könnte durch die Gradmessung noch der Nebenzweck erreicht werden, dass diese mit zur Vorbereitung tauglicher Personen für jenes Geschäft benutzt werden könnte.
